% A worked example of tikzflora: a short multipage document that uses every
% layout. Compile it with latexmk, which runs LuaLaTeX (latexmkrc) as
% often as the table of contents needs, or run lualatex on it twice.
%
% Every \florasetup and \florathispage below is commented with what it
% does; the full list of keys is in README.md and on
% https://flora.alpacasoft.dev/guide/.
%
% The text is in the public domain: Vitruvius, The Ten Books on
% Architecture, tr. Morris Hicky Morgan (1914), book IV, chapter 1; and
% Asa Gray, The Elements of Botany (1887), from the Project Gutenberg
% editions, abridged in places, with figure references and page markers
% left out.
\documentclass[11pt,twoside]{article}
% Two-sided pages with a wide outer margin: the flora grows there.
\usepackage[a4paper,inner=2cm,outer=5.6cm,top=4.2cm,bottom=3.5cm,
  headheight=14pt,headsep=18pt]{geometry}
\usepackage{lipsum}
\usepackage{fancyhdr}
\usepackage{tikzflora}

\pagestyle{fancy}
\fancyhf{}
\fancyhead[LE,RO]{\small\textsc{Florilegium}}
\fancyfoot[LE,RO]{\small\thepage}
\renewcommand\headrulewidth{0pt}

% The document's defaults. layout=split puts the flora in the outer margin:
% on the right of odd pages and the left of even ones. seed picks the
% garden; change it for a different one, or use seed=random. visitors=0.6
% invites a few bees, butterflies, birds and snails.
\florasetup{layout=split, seed=3, visitors=0.6}

\begin{document}

% --- Title page -----------------------------------------------------------
% A narrower, centred text block leaves room around the title. The flank
% layout grows a plant up each side of the page, their roots meeting at
% the foot round a crest. ornaments weighs which kind each place picks:
% here the sun crest weighs three times its usual, and the cul-de-lampe
% tailpiece too, for a layout that sets one (layout=title does).
\newgeometry{left=3.2cm,right=3.2cm,top=5cm,bottom=4cm}
\thispagestyle{empty}
\florathispage{layout=flank, ornaments={sun_crest=6, cul_de_lampe=3}}
\begin{floratitle}
\centering
{\Huge\scshape Florilegium\par}
\vspace{1.2em}
{\large Of leaves, after Vitruvius and Asa Gray\par}
\end{floratitle}
\vspace{6cm}
\begin{center}
\begin{minipage}{0.66\textwidth}
\small\itshape
The acanthus root, pressed down meanwhile though it was by the weight,
when springtime came round put forth leaves and stalks in the middle.
\par\smallskip\upshape\hfill Vitruvius, IV.1
\end{minipage}
\end{center}
\clearpage

% --- Contents -------------------------------------------------------------
% surround grows a wreath on all four sides of the text block.
\thispagestyle{empty}
\florathispage{layout=surround}
\tableofcontents
\clearpage
\restoregeometry

% --- Body -----------------------------------------------------------------
% The body uses the document default, layout=split.
\section{The acanthus}

It is related that the original discovery of this form of capital was as
follows. A free-born maiden of Corinth, just of marriageable age, was
attacked by an illness and passed away. After her burial, her nurse,
collecting a few little things which used to give the girl pleasure while
she was alive, put them in a basket, carried it to the tomb, and laid it on
top thereof, covering it with a roof-tile so that the things might last
longer in the open air. This basket happened to be placed just above the
root of an acanthus. The acanthus root, pressed down meanwhile though it was
by the weight, when springtime came round put forth leaves and stalks in the
middle, and the stalks, growing up along the sides of the basket, and
pressed out by the corners of the tile through the compulsion of its
weight, were forced to bend into volutes at the outer edges.

% A clearing: a gap taller than about three lines, as a figure would
% leave. The flora grows into it from the margin.
\vspace{5.2cm}
\begin{center}\small\itshape The basket and the tile, left to the margin.\end{center}

Just then Callimachus, whom the Athenians called \emph{katatexitechnos} for
the refinement and delicacy of his artistic work, passed by this tomb and
observed the basket with the tender young leaves growing round it.
Delighted with the novel style and form, he built some columns after that
pattern for the Corinthians, determined their symmetrical proportions, and
established from that time forth the rules to be followed in finished
works of the Corinthian order.

\section{The framework of the leaf}

The framework consists of wood,---a fibrous and tough material which runs
from the stem through the leaf-stalk, when there is one, in the form of
parallel threads or bundles of fibres; and in the blade these spread out in
a horizontal direction, to form the ribs and veins of the leaf. The stout
main branches of the framework are called the Ribs. When there is only one,
or a middle one decidedly larger than the rest, it is called the Midrib.
The smaller divisions are termed Veins; and their still smaller
subdivisions, Veinlets. The latter subdivide again and again, until they
become so fine that they are invisible to the naked eye. The fibres of
which they are composed are hollow; forming tubes by which the sap is
brought into the leaves and carried to every part.

Venation is the name of the mode of veining, that is, of the way in which
the veins are distributed in the blade. This is of two principal kinds;
namely, the parallel-veined, and the netted-veined.

% An indented quotation: the lines start further in, so on pages with the
% flora on the left it reaches into the indent.
\begin{quote}
In Netted-veined (also called Reticulated) leaves, the veins branch off
from the main rib or ribs, divide into finer and finer veinlets, and the
branches unite with each other to form meshes of network. That is, they
anastomose, as anatomists say of the veins and arteries of the body.
\end{quote}

In parallel-veined leaves, the whole framework consists of slender ribs or
veins, which run parallel with each other, or nearly so, from the base to
the point of the leaf,---not dividing and subdividing, nor forming meshes,
except by minute cross-veinlets. The leaf of any grass, or that of the Lily
of the Valley will furnish a good illustration. Such parallel veins
Linn\ae us called Nerves, and parallel-veined leaves are still commonly
called nerved leaves, while those of the other kind are said to be
veined,---terms which it is convenient to use, although these ``nerves''
and ``veins'' are all the same thing, and have no likeness to the nerves and
little to the veins of animals.

\section{Forms of leaves}

It is necessary to give names to the principal shapes, and to define them
rather precisely, since they afford easy marks for distinguishing species.
Beginning with the narrower and proceeding to the broadest forms, a leaf is
said to be

% Short lines in a list leave a ragged edge for the flora to fill.
\begin{description}
\item[Linear,] when narrow, several times longer than wide, and of the same
  breadth throughout.
\item[Lanceolate,] or Lance-shaped, when conspicuously longer than wide,
  and tapering upwards, or both upwards and downwards.
\item[Oblong,] when nearly twice or thrice as long as broad.
\item[Elliptical] is oblong with a flowing outline.
\item[Ovate,] when the outline is like a section of a hen's egg
  lengthwise, the broader end downward.
\item[Orbicular,] or Rotund, circular in outline.
\end{description}

% From here on: twice the lilies and umbels, no ghost blooms, a bird and a
% butterfly on most pages, and leaves in green inks. mix weights are 1 for
% the usual amount and 0 for none.
\florasetup{mix={lily=2, umbel=2, ghost=0, bird=1.5, butterfly=1.5},
  palette={leaves=green}}

\subsection{Margins}

As to particular outlines of simple leaves (and the same applies to their
separate parts), they are

\begin{itemize}
\item \emph{Entire}, when the margin is an even line, without teeth or
  notches.
\item \emph{Serrate}, or saw-toothed, when the margin is cut into sharp
  teeth, like those of a saw, and pointing forwards.
\item \emph{Dentate}, or toothed, when such teeth point outwards.
\item \emph{Crenate}, or scalloped, when the teeth are broad and rounded.
\item \emph{Repand}, undulate, or wavy.
\item \emph{Sinuate}, when the margin is more strongly sinuous.
\end{itemize}

\subsection{Lorem ipsum}
For comparison, a passage of the usual filler.

\lipsum[1-2]

% --- A centred page -------------------------------------------------------
% centered grows one plant in each side margin, on separate seeds. This
% page alone paints its blooms and hero in an ink picked at random; spin
% picks another without moving a single shape.
\clearpage
\newgeometry{left=4.6cm,right=4.6cm,top=4.2cm,bottom=3.5cm,headheight=14pt,headsep=18pt}
\florathispage{layout=centered, palette={blooms=random, hero=random}, spin=2}
\section{Tendrils}

Tendrils are branches of a very slender sort, like runners, not destined
like them for propagation, and therefore always destitute of buds or
leaves, being intended only for climbing. Simple tendrils are such as those
of Passion-flowers. Compound or branching tendrils are borne by the
Cucumber and Pumpkin, by the Grape-Vine, Virginia Creeper, etc.

A tendril commonly grows straight and outstretched until it reaches some
neighboring support, such as a stem, when its apex hooks around it to
secure a hold; then the whole tendril shortens itself by coiling up
spirally, and so draws the shoot of the growing plant nearer to the
supporting object. But the tendrils of the Virginia Creeper, as also the
shorter ones of the Japanese species, effect the object differently,
namely, by expanding the tips of the tendrils into a flat disk, with an
adhesive face. This is applied to the supporting object, and it adheres
firmly; then a shortening of the tendril and its branches by coiling
brings up the growing shoot close to the support. This is an adaptation
for climbing mural rocks or walls, or the trunks of trees, to which
ordinary tendrils are unable to cling. The Ivy and Poison Ivy attain the
same result by means of aerial rootlets.

\vspace{4.5cm}

\section{The umbel}

An Umbel, as in the Milkweed, is a sort of flower-cluster where the
pedicels all spring apparently from the same point, from the top of the
peduncle, so as to resemble, when spreading, the rays of an umbrella;
whence the name. Here the pedicels are sometimes called the Rays of the
umbel. And the bracts, when brought in this way into a cluster or circle,
form what is called an Involucre.

When the stalks which in the simple umbel are the pedicels of single
flowers themselves branch into an umbel, a Compound Umbel is formed. This
is the inflorescence of Caraway, Parsnip, and almost all of the great family
of Umbelliferous (umbel-bearing) plants.
\clearpage
\restoregeometry

% --- Closing --------------------------------------------------------------
% A short last page: the flora grows into the space below the text, here
% at a lighter density.
\florathispage{density=0.6}
\section{Colophon}
\raggedright
Set with \textsf{tikzflora}, whose margins are grown afresh on every page.

It would not exist without two sources: the r/LaTeX thread
\emph{How do you make these flowers in LaTeX?}
(\texttt{old.reddit.com/r/LaTeX/comments/1xixznw}) and Kevin Raymond's
\emph{florilegium} (\texttt{github.com/kevinraymond/florilegium}), whose grown
margins, animation and visitors this package follows. The ink colours are
Pilot iroshizuku ramps measured by madoromi
(\texttt{madoromi.moonwit.ch}); the growth follows Wong, Zongker and Salesin,
\emph{Computer-Generated Floral Ornament} (SIGGRAPH 1998).

The text is in the public domain: Vitruvius, \emph{The Ten Books on
Architecture}, translated by Morris Hicky Morgan (1914), and Asa Gray,
\emph{The Elements of Botany} (1887), both from Project Gutenberg,
abridged in places and with their figure references left out.

\end{document}
